\documentclass[12pt]{article}
\usepackage[utf8]{inputenc}
\usepackage[T1]{fontenc}
\usepackage[margin=1in]{geometry}
\usepackage{setspace}
\usepackage{xcolor}
\usepackage[colorlinks=true, urlcolor=blue!60!black, linkcolor=blue!60!black]{hyperref}
\setlength{\parindent}{0em}
\setlength{\parskip}{0.45em}
\renewcommand{\baselinestretch}{1.0}
\begin{document}
\begin{center}
{\Large \textit{Medell\'{i}n v. Texas}, 552 U.S. 491 (2008)\par}
\vspace{0.3em}
{\normalsize Case Brief --- One-Page Edition\par}
\end{center}

\textbf{Memory Jogger.} \textit{Avena}, the President, and non-self-executing treaties (\href{https://supreme.justia.com/cases/federal/us/552/491/}{552 U.S.\ 491}).

\textbf{Facts.} Jos\'{e} Ernesto Medell\'{i}n, a Mexican national on Texas's death row, was never told of his Vienna Convention right to consular notice. The ICJ's \textit{Avena} judgment required ``review and reconsideration'' of his case, and President Bush directed state courts to comply.

\textbf{Procedural History.} Texas courts dismissed his second habeas application as an abuse of the writ.

\textbf{Votes.} Affirmed, 6--3 (Roberts for the Court; Stevens concurring in the judgment; Breyer, Souter, and Ginsburg dissenting).

\textbf{Holding.} Neither \textit{Avena} nor Article 94 of the U.N.\ Charter is self-executing domestic law, and the President, in \textit{Youngstown}'s Category Three, cannot unilaterally make a non-self-executing treaty preempt state law.

\textbf{Reasoning.} ``Undertakes to comply'' commits the political branches, and the Charter's remedy is diplomatic. Converting a treaty into domestic law conflicts with the ratifying Senate's understanding and belongs to Congress; claims-settlement cases do not apply. Breyer urged a practical test and warned that some 70 treaties were at risk.

\textbf{Examples.} \emph{Same-Side:} a directive enforcing a tribunal ruling. \emph{Opposite-Side:} implementing legislation like the ICSID statute. \emph{Fence-Sitter:} \S\ 1983 suits for denial of consular notice.

\textbf{Critique.} One state can put the Nation in breach; critics see a presumption against self-execution.

\textbf{Key Quotation.} ``In sum, while treaties `may comprise international commitments . . . they are not domestic law unless Congress has either enacted implementing statutes or the treaty itself conveys an intention that it be ``self-executing'' and is ratified on these terms.' '' --- Chief Justice Roberts, 552 U.S.\ at 505.

\end{document}
